	\section{Spinons}

The strong coupling between plasmons and magnetic excitations we describe here is a general property of a plasmon interacting with spin excitations. To showcase this, we now consider a completely different kind of spin excitation: the spinon. Specifically, we consider a layered spin-1/2 antiferromagnetic Heisenberg model on the Kagome lattice, at half filling, which has been proposed to be realized by Herbertsmithite \cite{normanColloquiumHerbertsmithiteSearch2016}. Previous studies suggest that the ground state is a U(1)-Drac spin-liquid state, which elementary excitation are spinons \cite{ranProjectedWaveFunctionStudySpin$12007}. Since the Kagome lattice has broken inversion symmetry along its bonds, this makes it an excellent candidate to demonstrate how different spin excitations can couple to plasmons. We assume the layers to be effectively decoupled, as is typically the case in spin-liquid candidates.


\begin{figure}
	\centering
	\includegraphics{dispersion_spinons}
	\caption{The spinon dispersion, relative to the bottom of the first band. The colorscale is the spinon 2-DOS. Shown on the right is the spinon 2-DOS for $\bm k=\Gamma$. \label{fig:spinons}}
\end{figure}

We first describe a single layer, where the dynamics are given by a nearest-neighbhor Heisenberg Hamiltonian 
\begin{equation}
	H = J \sum_{\langle i,j\rangle} \hat{\bm S}_i \cdot \hat{\bm S}_j,
\end{equation}
with antiferromagnetic exchange $J>0$.
Next, we introduce the fermionic spinon operators such that
\begin{equation}
	\hat{\bm S}_i = \sum_{\alpha\beta} \hat f^\dagger_{i\alpha} \bm\sigma_{\alpha\beta}\hat f_{i\beta}, \label{eq:S-spinons}
\end{equation}
together with the constraint that $\sum_\alpha \hat f_{i\alpha}^\dagger \hat f_{i\alpha}=1$. here $\bm\sigma_{\alpha\beta}$ is a vector of Pauli matrices. We then obtain the interacting spin Hamiltonian
\begin{equation}
	H = \frac{J}{4}\sum_{ij}\sum_{\alpha\beta}  \left[2 \hat f_{i\alpha}^\dagger \hat f_{i\beta} \hat f_{j\beta}^\dagger \hat f_{j\alpha} - \hat n_{i\alpha}\hat n_{j\beta} \right],
\end{equation}
where $\hat n_{i\alpha}\equiv\hat f_{i\alpha}^\dagger \hat f_{i\alpha}$. 
This Hamiltonian can then be treated in a mean field approach, by defining a spinon hopping amplitude $\chi_{ij}\equiv\sum_\alpha\langle f_{i\alpha}^\dagger f_{j\alpha}\rangle$. It was shown by \textcite{ranProjectedWaveFunctionStudySpin$12007} that the energetically favorable state is a U(1)-Dirac spin-liquid state, described by zero flux through the triangles and $\pi$ flex through the hexagons. Here flux through a plaquette refers to the phase of $\chi_{ij}$ around the plaquette, since the phase of $\chi_{ij}$ can be viewed as the spatial components of a compact U(1) gauge field.

The U(1)-Dirac spin liquid can now be written as a mean-field Hamiltonian
\begin{equation}
	H = J \sum_{\langle ij\rangle} \chi_{ij} \hat f_{j\alpha}^\dagger \hat f_{i\alpha} + \hc,
\end{equation}
where the $ \chi_{ij}$ all have the same magnitude, with signs such that they produce zero flux through the triangles and $\pi$ flex through the hexagons, as indicated in Fig.??? Note that the unit cell has to doubled in order to fulfill this constraint. We now perform the Fourier transformation $\hat f_i=\sum_{\bm k}e^{i\bm r_i\cdot\bm k} \hat f_{\bm k}$ to find 
\begin{equation}
	H = J\sum_{\bm k}\sum_{nn'}\sum_\alpha \hat f_{\bm k;n\alpha}^\dagger \chi_{\bm k}^{nn'}\hat f_{\bm k;n'\alpha},
\end{equation}
where $\chi_{\bm k}^{nn'}=\sum_{\bm r_{ij}}e^{-i\bm r_{ij}\cdot \bm k} \chi_{ij}$, which we give explicitly in \cref{app:spinons}. To diagonalize this mean-field Hamiltonian, we introduce the normal modes
$\hat g_{\bm k;\mu\alpha}^\dagger\equiv \sum_{n} \mathcal U_{\bm k;\mu n}\hat f_{\bm k;n\alpha}^\dagger $, defined such that $\mathcal U_{\bm k}^\dagger \chi_{\bm k}\mathcal U_{\bm k}=\diag[\epsilon_1,\dots,\epsilon_6]$.

Having obtained the mean-field Hamiltonian, we show the spinon dispersion in \cref{fig:spinons}. We first observe the existence of a flat band at high energy, which is double degenerate. Secondly, there are two Dirac points, where the spinon dispersion is linear. The Fermi level corresponding to one spinon per site is is located at the Dirac point, such that the lower three bands are occupied.

\subsection{Spinon-plasmon coupling}
Analagous to the derivation for the magnon-plasmon coupling, we now write the polarization Hamiltonian \cref{eq:Hc-AFM-final} by introducing the spinon operators defined in \cref{eq:S-spinons}
\begin{equation}
	H = -\sum_{ij}\sum_{\alpha\beta\mu\nu} (\hat{\bm E}_i\cdot\bm\pi_{ij}) (\hat f_{i\alpha}^\dagger \bm\sigma_{\alpha\beta} \hat f_{i\beta}) \cdot(f_{j\mu}^\dagger\bm \sigma_{\mu\nu}f_{j\nu}).
\end{equation}
Note that in the Kagome lattice, the polarization vectors point into the triangles, as indicated in Fig.???.
Since the coupling Hamiltonian we obtain is quartic in spinon operators, we again make a mean-field assumption, and find 
\begin{equation}
	H_c = -\sum_{\alpha}\sum_{ij} (\hat{\bm E}(\bm r_i)\cdot \bm \pi_{ij}) \chi_{ij} \left(f_{j\alpha}^\dagger f_{i\alpha} + \hc \right)
\end{equation}
which we Fourier transform and write in the normal basis as
\begin{equation}
	H_c = -i\sum_\alpha\sum_{\mu\nu}\sum_{\bm k\bm q} (\hat\phi_{\bm q}-\hat\phi_{-\bm q}) ( \bm E_{\bm q}\cdot\bm{\mathcal V}_{\bm k + \bm q}^{\mu\nu})( g_{\bm k+\bm q;\mu\alpha}^\dagger g_{\bm k;\nu\alpha} + \hc)
\end{equation}
where we have defined
\begin{equation}
	{\bm{\mathcal V}}_{\bm k + \bm q}^{\mu\nu} = \sum_{nn'}\mathcal U_{\bm k + \bm q;n\mu}^*\bar{\bm\pi}_{\bm k+\bm q}^{nn'} \mathcal U_{\bm k;n'\nu}
\end{equation}
as the interaction vertex in the diagonalized normal basis, with
\begin{equation}
	\bar{\bm\pi}_{\bm k}^{nn'} = \sum_{\bm r_{ij}}e^{-i\bm r_{ij}\cdot\bm k} \chi_{ij}^{nn'}\bm\pi_{ij}^{nn'}
\end{equation}
the polarization vector in mean field. We give $\bar{\bm\pi}_{\bm k}^{nn'}$ in \cref{app:spinons}.

From the coupling Hamiltonian we can now write down the plasmon self-energy for $N_z$ Kagome layers directly
\begin{multline}
	\Sigma_{\bm q}(\epsilon) = \frac{2N_z}{N}\sum_{\bm k} \sum_{\mu\nu} \frac{|\bm E_{\bm q}\cdot {\bm{\mathcal V}}_{\bm k + \bm q}^{\mu\nu}|^2}{\epsilon + \ii 0^+ + \epsilon_{\bm k;\mu} - \epsilon_{\bm k+\bm q;\nu}}\times \\ 
	\left(n_F(\epsilon_{\bm k;\mu}) - n_F(\epsilon_{\bm k+\bm q;\nu})\right)
\end{multline}
where $n_F$ is the Fermi-Dirac distribution. At half filling and zero temperature, the lower 3 bands are filled, and thus the sum reduces to $\mu\in\{1,2,3\}$ and $\nu\in\{4,5,6\}$. Note that since we have assumed the interlayer exchange coupling to be negligible, we obtain a simple linear in $N_z$ scaling of the self-energy. The factor of $2$ follows from summing over the spin degree of freedom $\alpha$.

Analogous to the magnon case, we observe that the imaginary part of the plasmon self-energy,
\begin{equation}
	\Im[\Sigma_{\bm q}(\epsilon)] = \frac{2N_z}{N}\sum_{\bm k}\sum_{\substack{\mu\leq3\\\nu>3}} {|\bm E_{\bm q}\cdot\bm{\mathcal V}_{\bm k+\bm q}^{\mu\nu}|^2}\delta({\varepsilon -\varepsilon^\mu_{\bm k} + \varepsilon^{\nu}_{\bm k+\bm q}}). 
\end{equation}
is directly related to the spin-zero two-spinon density of states (2-DOS) defined as
\begin{equation}
	D_{\bm q}(\epsilon)= \frac1N \sum_{\substack{\mu\leq3\\\nu>3}}\sum_{\bm k} \delta (\epsilon - \varepsilon^\mu_k - \varepsilon^{\nu}_{\bm + \bm q}).
\end{equation}
The 2-DOS therefore also plays an important role in determining the plasmon self-energy. 

\subsection{Results}
We now evaluate the plasmon self-energy numerically, and show the resulting plasmon spectral density in  \cref{fig:spinon-spectral}. The upper edge of the spinon continuum is indicated by the horizontal white line. 
We first observe the existence of multiple points in the spectrum where the plasmon gets expelled, in contrast to the magnon-plasmon coupling. We attribute this to the more complicated spectrum of the spinon continuum, which we investigate in more detail by calculating the spinon 2-DOS, which we show in \cref{fig:spinons}. We observe the existence of several peaks in the 2-DOS at $\bm q=\Gamma$, which can be attributed to the flat band transitions between the lower filled bands and the upper unfilled bands. In particular, we expect six of such transitions, since the upper band is degenerate, and we can indeed distuingish six peaks in the spinon 2-DOS. Note that even though transitions between the upper and lower branch of the Dirac point are allowed, these have vanishing support once integrated over the Brillouin zone, and thus there is no noticable peak of the 2DOS at zero energy. 

However, the structure of the spinon 2DOS does not explain the multitude of fine-grained features we observe in \cref{fig:spinon-spectral}. We note however that in order for the plasmon to scatter a spinon from a filled band $\mu$ to an unfilled band $\nu$, the matrix element $\bm E_{\bm q}\cdot \mathcal V_{\bm k + \bm q}^{\mu\nu}$ has to be non-zero. In particular, $\mathcal V_{\bm k + \bm q}^{\mu\nu}$ is dependent on the microscopic details, and the fine-grained features are thus not generic feature of plasmon-spinon interactions.

Finally, we note that upon exciting the spinon continuum the plasmon dispersion remains modified, since the real part of the self-energy remains finite even outside the continuum. The imaginary part of the self energy is zero however, and thus the lifetime of the plasmon is unaffected.

In conclusion, we have shown that it the spinon continuum couples to plasmons, expelling and decaying the plasmon. The strength of the coupling is however considerably weaker than the in the magnon case. This we attribute to the fact that for the magnons, all the magnon


\begin{figure}
	\includegraphics[width=\columnwidth]{spinon_spectral_Nz_1000}
	\caption{The plasmon spectral function for a $N_z=1000$ layered kagome antiferromagnet. The horizontal solid line indicates the bandwith of the spinon spectrum. The dashed white line is the bare plasmon dispersion. \label{fig:spinon-spectral}}
\end{figure}
\subsection{Extra}
\label{app:spinons}
\begin{widetext}
	\begin{equation}
		\chi_{\bm k}^{nn'} = \begin{pmatrix}
			0 & e^{i\bm \tau_1 \cdot \bm k} & e^{i\bm \tau_2 \cdot \bm k} - e^{i\bm \tau_5 \cdot \bm k} & 0 & e^{i\bm \tau_4 \cdot \bm k} & 0 \\
			e^{-i\bm \tau_1 \cdot \bm k} & 0 & e^{i\bm \tau_3 \cdot \bm k} & -e^{i\bm \tau_1 \cdot \bm k} & 0 & -e^{i\bm \tau_6 \cdot \bm k} \\ 
			e^{-i\bm \tau_2 \cdot \bm k} - e^{-i\bm \tau_5 \cdot \bm k} & e^{i\bm \tau_3 \cdot \bm k} & 0 & 0 & -e^{-i\bm \tau_3 \cdot \bm k} & 0 \\
			0 & -e^{-i\bm \tau_1 \cdot \bm k} & 0 & 0 & e^{i\bm \tau_1 \cdot \bm k} & e^{i\bm \tau_2 \cdot \bm k} +e^{i\bm \tau_5 \cdot \bm k} \\ 
			e^{-i\bm \tau_4 \cdot \bm k} & 0 & -e^{-i\bm \tau_3 \cdot \bm k} & e^{-i\bm \tau_1 \cdot \bm k} & 0 & e^{i\bm \tau_3 \cdot \bm k} \\ 
			0 & -e^{-i\bm \tau_6 \cdot \bm k} & 0 & e^{-i\bm \tau_2 \cdot \bm k} + e^{-i\bm \tau_5 \cdot \bm k} & e^{-i\bm \tau_3 \cdot \bm k} & 0
		\end{pmatrix}
	\end{equation}
	The polarization in mean field can then be found as 
	\begin{equation}
		\bar{\bm \pi}^{nn'}_{\bm k} = g\begin{pmatrix}
			0 & \bm\nu_1 e^{i\bm \tau_1 \cdot \bm k} & \bm\nu_2 e^{i\bm \tau_2 \cdot \bm k} + \bm\nu_5 e^{i\bm \tau_5 \cdot \bm k} & 0 & \bm\nu_4 e^{i\bm \tau_4 \cdot \bm k} & 0 \\
			\bm\nu_1 e^{-i\bm \tau_1 \cdot \bm k} & 0 & \bm\nu_3 e^{i\bm \tau_3 \cdot \bm k} & \bm\nu_1 e^{i\bm \tau_1 \cdot \bm k} & 0 & \bm\nu_6 e^{i\bm \tau_6 \cdot \bm k} \\ 
			\bm\nu_2 e^{-i\bm \tau_2 \cdot \bm k} - \bm\nu_5 +^{-i\bm \tau_5 \cdot \bm k} & \bm\nu_3 e^{i\bm \tau_3 \cdot \bm k} & 0 & 0 & -\bm\nu_3 e^{-i\bm \tau_3 \cdot \bm k} & 0 \\
			0 & \bm\nu_1 e^{-i\bm \tau_1 \cdot \bm k} & 0 & 0 & \bm\nu_1  e^{i\bm \tau_1 \cdot \bm k} & \bm\nu_2 e^{i\bm \tau_2 \cdot \bm k} + \bm\nu_5 e^{i\bm \tau_5 \cdot \bm k} \\ 
			\bm\nu_4 e^{-i\bm \tau_4 \cdot \bm k} & 0 & -\bm\nu_3 e^{-i\bm \tau_3 \cdot \bm k} & \bm\nu_1 e^{-i\bm \tau_1 \cdot \bm k} & 0 & \bm\nu_3 e^{i\bm \tau_3 \cdot \bm k} \\ 
			0 & \bm\nu_6 e^{-i\bm \tau_6 \cdot \bm k} & 0 & \bm\nu_2 e^{-i\bm \tau_2 \cdot \bm k} +\bm\nu_5 e^{-i\bm \tau_5 \cdot \bm k} & \bm\nu_3 e^{-i\bm \tau_3 \cdot \bm k} & 0
		\end{pmatrix}
	\end{equation}
	
	where $\bm\tau_i= a[(1,0), (\frac12, \frac{\sqrt3}{2}),(-\frac12, \frac{\sqrt3}{2}), (-1,0), (-\frac12,- \frac{\sqrt3}{2}),(\frac12, - \frac{\sqrt3}{2})  ]$ 
	and $\bm \nu_i = \mathcal R_{\pi/2}\bm \tau_i$, where $\mathcal R_{\pi/2}$ rotates at angle $\pi/2$. Additionally, $g=\cos\frac\pi3 \left(\frac{t}{U}\right)^3 $. Note that the factor $\cos\frac\pi3 $ follows from the projection of the ???
	%	where $\mathcal R_{\pi/2}$ rotates at an angle.	
\end{widetext}


\section{Hematene}
We use the crystal structure as determined by \textcite{gonzalezHematene2DMagnetic2019}. Hematene has a trigonal unit cell, described by $\bm a_1=a\left(\frac12,- \frac{\sqrt{3}}{2}\right)$ and $\bm a_2=a\left(\frac12, \frac{\sqrt{3}}{2}\right)$, with four iron atoms, which we label $A$ through $D$ in this work, and six oxygen atoms, numbered $1$ through $6$. The positions of the atoms is given in \cref{tab:hematene}. 
\begin{table}[]
	\centering
	\begin{tabular}{ccc}
		Atom & Label & Position  \\
		\ch{Fe} & $A$ &  $\eta\hat{\bm z}$ \\ 
		\ch{Fe} & $B$ & $-\eta\hat{\bm z}$\\
		\ch{Fe} & $C$ & $\frac13 \bm a_1 + \frac23 \bm a_2 $\\
		\ch{Fe} & $D$ & $\frac23 \bm a_1 + \frac13 \bm a_2$ \\
		\ch{O} & 1 & $(\frac23-\delta)\bm a_2+\zeta\hat{\bm z}$ \\
		\ch{O} & 2 & $(\frac13+\delta)\bm a_1+(\frac13+\delta)\bm a_2+\zeta\hat{\bm z}$ \\
		\ch{O} & 3 & $(\frac231\delta)\bm a_1+\zeta\hat{\bm z}$ \\
		\ch{O} & 4 & $(\frac13+\delta)\bm a_2-\zeta\hat{\bm z}$\\
		\ch{O} & 5 & $(\frac13+\delta)\bm a_1-\zeta\hat{\bm z}$\\
		\ch{O} & 6 & $(\frac23-\delta)\bm a_1+(\frac23-\delta)\bm a_2-\zeta\hat{\bm z}$
	\end{tabular}
	\caption{Here, $\eta=\SI{1.56}{\AA}$, $\zeta=\SI{1.13}{\AA}$ and $\delta=0.0132$, $\bm a_1=a\left(\frac12,- \frac{\sqrt{3}}{2}\right)$ and $\bm a_2=a\left(\frac12, \frac{\sqrt{3}}{2}\right)$, with $a=\SI{5.09}{\AA}$.}
	\label{tab:hematene}
\end{table}

We use in this work the proposed antiferromagnetic ordering by \textcite{gonzalezHematene2DMagnetic2019}, where the \ch{Fe}-atoms $A$ and $B$ are aligned, as well as $C$ and $D$, but $A,B$ and $C,D$ are anti-aligned. In addition, there are three nearest-neighbhor exchange interactions: (i) between $A$ and $B$, parametrized by $J_{AB}$, (ii) between $C$ and $D$, parametrized by $J_{CD}$ and (iii) between $A,B$ and $C,D$, parametrized by $J'$. \textcite{gonzalezHematene2DMagnetic2019} find all three exchange interactions to be antiferromagnetic, with values $J_{AB}=\SI{7.2}{meV}$, $J_{CD}=\SI{3.2}{meV}$ and $J'=\SI{15.3}{meV}$. In addition, we consider an on-site magnetic anisotropy, parameterized by $K$, such that the exchange Hamiltonian becomes
\begin{equation}
	H = \sum_{ij} J_{ij} S_i S_j + \frac{K}{2} \sum_i S_z^2.
\end{equation}


We can now introduce the Holstein-Primakoff operators and find after Fourier transforming
\begin{equation}
	H = \sum_{\bm k} \psi_{\bm k}^\dagger \begin{pmatrix}
		A_{\bm k} & B_{\bm k} \\
		B_{\bm k}^\dagger & A_{\bm k}
	\end{pmatrix} \psi_{\bm k},
\end{equation}
where 
\begin{equation}
	A_{\bm k} \equiv  \begin{pmatrix}
		E_{AB} & -J_{AB}\gamma_{AB}(\bm k) & 0 & 0 \\
		-J_{AB}\gamma_{AB}^*(\bm k) & E_{AB} & 0 & 0 \\
		0 & 0 & E_{CD} & -J_{CD}\gamma_{CD}(\bm k) \\
		0 & 0 & -J_{CD}\gamma_{CD}^*(\bm k) & E_{CD}
	\end{pmatrix}
\end{equation}
and 
\begin{equation}
	B_{\bm k} \equiv \begin{pmatrix}
		0 & 0 & J'\gamma_{AC}(\bm k) & J'\gamma_{AD}(\bm k) \\
		0 & 0 & J'\gamma_{BC}(\bm k) &  J'\gamma_{BD}(\bm k) \\
		J'\gamma_{AC}^*(\bm k) & J'\gamma_{BC}^*(\bm k) & 0 & 0\\
		J'\gamma_{AD}^*(\bm k) & J'\gamma_{BD}^*(\bm k) & 0 & 0
	\end{pmatrix}.
\end{equation}
Additionally, 
\begin{equation}
	\psi_{\bm k} = \begin{pmatrix}
		a_{A\bm k} \\ a_{B\bm k} \\ b_{C\bm k} \\ b_{D\bm k} \\ a_{A-\bm k}^\dagger \\ a_{B-\bm k}^\dagger \\ b_{C-\bm k}^\dagger \\ b_{D-\bm k}^\dagger
	\end{pmatrix}.
\end{equation}
Here $E_{AB} \equiv K + J_{AB} + 6J'$, $E_{CD} \equiv K +  3J_{CD} + 6J'$ and $\gamma_{\eta\mu}\equiv \sum_{\eta\mu}e^{i\bm k\cdot \bm r_{\eta\mu}}$ is a sum over the nearest neighbhors between $\eta,\mu$.

Following \textcite{zhuSpinorbitCouplingElectronic2014}, we determine the direction of the polarization operator connecting two \ch{Fe}-atoms by forming a loop, which will consist of two \ch{Fe}-atoms and one \ch{O}-atom. We can then obtain the directions of the polarization operators, as given in \cref{tab:polarization-hematene}.
\begin{table}[]
	\centering
	\begin{tabular}{ccc}
		Atom pair & $\hat{\bm P_{ij}}$ & $\bm r_{ij}$ \\
		$C\rightarrow A$ & (0.196, 0.399, 0.896) & (0.5, -0.289, 1.56)\\ 
		$C\rightarrow A$ & (0.247, -0.369, 0.896) & (-0.5, -0.289, 1.56)\\ 
		$C\rightarrow A$ & (-0.443, -0.029, 0.896) & (0.0, 0.577, 1.56)\\ 
		$C\rightarrow B$ & (0.196, 0.399, 0.896) & (0.0, 0.577, -1.56)\\ 
		$C\rightarrow B$ & (0.247, -0.369, 0.896) & (0.5, -0.289, -1.56)\\ 
		$C\rightarrow B$ & (-0.443, -0.029, 0.896) & (-0.5, -0.289, -1.56)\\ 
		$D\rightarrow B$ & (0.196, 0.399, 0.896) & (0.0, -0.577, -1.56)\\ 
		$D\rightarrow B$ & (0.247, -0.369, 0.896) & (-0.5, 0.289, -1.56)\\ 
		$D\rightarrow B$ & (-0.443, -0.029, 0.896) & (0.5, 0.289, -1.56)\\ 
		$D\rightarrow A$ & (0.196, 0.399, 0.896) & (0.0, -0.577, 1.56)\\ 
		$D\rightarrow A$ & (0.247, -0.369, 0.896) & (0.5, 0.289, 1.56)\\ 
		$D\rightarrow A$ & (-0.443, -0.029, 0.896) & (-0.5, 0.289, 1.56)\\ 
	\end{tabular}
	\caption{For the \ch{Fe}-atom pairs, the polarization vector $\hat{\bm P_{ij}}$ and corresponding vector $\bm r_{ij}$ connecting the pair. $\hat{\bm P_{ij}}$ is given as a unit vector and $\bm r_{ij}$ has units \si{\AA}.}
	\label{tab:polarization-hematene}
\end{table}

\section{Antiferromagnet on square lattice}
We will first demonstrate the induced lifetime broadening by considering a square antiferromagnet in the classical limit, where the elementary spin excitations are magnons. In addition, we break the bond inversion symmetry by placing a non-magnetic atoms as shown in Fig.??. This allows us to describe all the essential features of the lifetime broadening, before moving on to more realistic setups describing excitations in Hematene and ??? in Secs...




The spin degree of freedom is described by the Heisenberg Hamiltonian on a square lattice, 
\begin{equation}
	H = \sum_{ij} J_{ij}\bm S_i \cdot \bm S_j 
\end{equation}
where $J_{ij}>0$ for the nearest neighor spins and zero otherwise. The ground state is therefore an antiferromagnetic ordering of the spins. In the limit of large $S$ the elementary magnetic excications of this system are magnons, which we introduce using the Holstein-Primakoff operators $S_{A;i}^+ \approx \sqrt{2S} a_i $, $S_{A;i}^z\approx S- a_i^\dagger a_i$, $S_{B;i}^+ \approx\sqrt{2S}  b_i^\dagger $ and $S_{B;i}^z\approx -S+ b_i^\dagger b_i$. After Fourier transformation, the Hamiltonian $H=\sum_{\bm k} J(a_{\bm k}a_{\bm k}+a_{\bm k}a_{\bm k})+J_{\bm k}(a_{\bm k}b_{-\bm k} + a_{\bm k}^\dagger b_{-\bm k}^\dagger)$ can be diagonalized through the Boguliobuv transformation $a_{\bm k} = u_{\bm k} \alpha_{\bm k} - v_{\bm k} \beta_{-\bm k}^\dagger$ and $b^\dagger_{-\bm k} = -v_{\bm k} \alpha_{\bm k} + u_{\bm k} \beta_{-\bm k}^\dagger$ to find
\begin{equation}
	H=\sum_{\bm k} \epsilon_{\bm k} \left( \alpha_{\bm k}^\dagger \alpha_{\bm k} + \beta_{\bm k}^\dagger\beta_{\bm k} \right)
\end{equation}
Here, 
\begin{equation}
	\epsilon_{\bm k} =  \sqrt{}\label{eq:dispersion-afm-square}
\end{equation}
is the antiferromagnetic dispersion relation and 
\begin{equation}
	u_{\bm k} = ??;\quad v_{\bm k}=??. \label{eq:bog-factors}
\end{equation}

We can now introduce the coupling between the plasmon and magnons by noting that a plasmon generates an electric field given by
\begin{equation}
	\bm E_i^\alpha = \frac{1}{\sqrt{N_p}}\sum_{\bm q} i \bm E_{\bm q} e^{i\bm q \cdot \bm r_i} \hat\phi_{\bm q} + \hc
\end{equation}
with $\hat\phi_{\bm q}$ the plasmon annihilation operator and $\bm E_{\bm q}$ is a mode function, which captures the details of the plasmonic dispersion. In what follows, we will assume the plasmons to exist in monolayer graphene, for which $\bm E_{\bm q}$ is known \cite{ferreira}. We repeat expressions for $\bm E_{\bm q}$ in the appendix ??. In general, $\bm E_{\bm q}$ can be separated into two parts, $\bm E\equiv E_{\bm q}^\parallel \hat{\bm q} + E_{\bm q}^z \hat{\bm z}$, where $\hat{\bm q}$ is the plasmon propagation direction. In what follows, we will assume that the plasmon length scale is much larger then the magnon length scale, i.e., the plasmon wavelength is much larger than the magnon wavelengths.  \pg{explain further here that we assume $E_i$ does not change too much in space?}



This allows to write the coupling Hamiltonian as
\begin{equation}
	H_c = -i\frac{1}{\sqrt{N_p}}\sum_{ij}\sum_{\bm q}\sum_{\alpha\beta} (\bm\pi_{ij}\cdot \bm F_{\bm q}) e^{i\bm q\cdot\bm r_i} \hat\phi_{\bm q} \hat{\bm S_i} \hat{\bm S_j} + \hc,
\end{equation}
which thus couples the plasmonic and spin operators. In this work, we consider antiferromagnets, which universally allow for magnon pair creation. As such, we only consider the isotropic part $\bm\pi_{ij}$. We now expand the spin operators in Holstein-Primakoff operators, Fourier transform and apply the Boguliobov transformation, such that we obtain 
\begin{multline}
	H_c = \frac{-i}{\sqrt{N_p}}\sum_{\bm k \bm q} (\bm F_{\bm q}\cdot \bm \pi_{\bm q}) (\hat\phi_{\bm q} + \hat\phi_{\bm q}^\dagger) \times \\
	(u_{\bm q}u_{\bm k-\bm q}\alpha_{\bm k} \beta_{\bm q - \bm k} + v_{\bm q}v_{\bm k-\bm q} \alpha_{\bm k}^\dagger \beta_{\bm q - \bm k}^\dagger 
	-u_{\bm q}v_{\bm k-\bm q}\alpha_{\bm k}^\dagger \alpha_{\bm q - \bm k} - v_{\bm q}u_{\bm k-\bm q} \beta_ {\bm k}^\dagger \beta_{\bm q - \bm k}) + \hc 
\end{multline}
The coupling Hamiltonian is first order in plasmon operators, and second order in magnon operators. Our goal is to derive the plasmonic self-energies at zero temperature, of which there are two kinds, depicted in \cref{fig:feynmann}: (1) forward and (2) circle diagrams. Since magnons are non-conserved bosonic particles, at zero temperature no magnons are present and thus the circle diagram does not contribute. 

We can now directly write down the plasmon self-energy from the forward diagram, which at zero temperature is
\begin{equation}
	\Sigma_{\bm q} = \frac{N_z}{N}\sum_{\bm k}(\bm F_{\bm q}\cdot \bm \pi_{\bm q})^2 \frac{\mathcal V_{q,k}^2}{\varepsilon + \ii 0^+ -\varepsilon_k - \varepsilon_{q-k}}. \label{eq:self-energy}
\end{equation}
Here $ \mathcal V_{q,k}$ is the interaction vertex in the Boguliobov basis, given by $ \mathcal V_{q,k}\equiv u_k u_{k-q} +v_k v_{k-q}$. $N_z$ is the number of layers of the antiferromagnet, which enters the interaction vertex as $\sqrt{N_z}$, and since the interaction vertex enters twice in the self-energy, we obtain a linear in $N_z$ scaling.

This corresponds to a single plasmon coupling to $N_z$ layers (provided the electric field generated by the plasmons extends along all $N_z$ layers). In the presence of an interlayer exchange coupling---either ferromagnetic or antiferromagnetic---the layers will couple and create $N_z$ modes, with the $n$th mode at energy $J_z S \pi^2 (n/N_z)^2$. For sufficiently strong $J_z$, the higher modes are therefore decoupled from the lowest, Kittel-like, mode and we can approximate the coupling of the plasmon to the stack players by coupling to the Kittel mode. \pg{Depending on the number of layers, the Kittel mode might not be sufficiently separated.}


Alternatively, we propose to create a stack of decoupled antiferromagnetic layers---similar to how the inelastic neutron scattering crosssection is increased for magnetic van der Waals materials by creating stacks of single monolayers. In this setup, a single interaction vertex would scale as $N_z$ and thus the plasmon self-energy would scale as $N_z^2$.

In the intermediate regime, where $J_z$ is neither strong enough to fully gap out the second mode nor weak enough for the layers to be treated indeppendenlty, we expect that the interaction vertex scaling is between $\sqrt{N_z}$ and linear and thus the self energy would scale at least linearly, and at most quadratically. In what follows, we will assume that the modes are coupled strong enough that we couple only to the Kittel mode and thus self-energy in \cref{eq:self-energy} scales as $\propto N_z$.



\pg{Since we assume that $q\ll k$, where $k$ is the magnon wavelength, we can further approximate this as $\mathcal V_{q,k}\approx u_k^2 + v_k^2$, which is otherwise known as the antiferromagnetic exchange enhancement factor.}






We are now in the position to evaluate it numerically. To achieve this, we sum \cref{eq:self-energy} over the full magnon Brillouin zone, where $\epsilon_k$ is given by \cref{eq:dispersion-afm-square} and $u_k$ and $v_k$ by \cref{eq:bog-factors}. Additionally, the placing of a non-magnetic atom introduces a polarization as $F_{ij}=g \bm{\hat{z}}$, where we have parameterized the polarization with the phenomenological parameter $g$. After Fourier transformation, we obtain $\bm\pi_{\bm q}=2g \left(\cos (q_x a) + \cos (q_y a)\right) \hat{\bm{z}}$. 

\begin{figure}
	\centering
	\includegraphics[width=\columnwidth]{plasmon_spectral}
	\caption{The plasmon spectral density as a function of wave vector $q$ and energy $\hbar\omega$. The horizontal lines indicate the upper and lower bounds of the two-magnon density of states. Results shows are for $N_z=100$ layers.}
	\label{fig:plasmon-spectral-density}
\end{figure}

In \cref{fig:plasmon-spectral-density} we show the resulting plasmon spectral density,
\begin{equation}
	A_p \equiv \frac{1}{\pi} \Im G = \Gamma \frac{1}{(\hbar\omega-\hbar\omega_{p,\bm q})^2+\Gamma^2}
\end{equation}
as a function of plasmon wavelength $q$ and frequency. We consider a stack of $N_z=???$ layers, as well as a intrinsic plasmon damping of $\Gamma_i=\gamma\hbar\omega$ with a quality factor $\gamma=10^{-3}$. 

We first observe that for low plasmon frequencies, there is no notable self-energy and thus the plasmon is well-defined. Upon entering the two-magnon continuum, the plasmon spectral density gets smeared out, because of the induced self-energy. 

This induced plasmon lifetime gradually reduces, and in fact drops below the intrinsic lifetime for a specific range of plasmon frequencies. At higher frequencies, there is a characteristic spike, and the lifetime start to rise sharply again. 

\begin{figure}
	\centering
	\includegraphics[width=\columnwidth]{dispersioncut}
	\caption{The imaginary (solid) and real (dashed) part of the on-shell plasmon self-energy $\Gamma$ induced by the magnon-plasmon coupling. Also shown are the plasmon energies and the quality factor. Here we have assumed an intrinsic lifetime of $\gamma_i=10^{-2}$.  }
	\label{fig:enter-label}
\end{figure}

These features can be most easily understood by noting that the imaginary part of the plasmon self-energy can be written as 
\begin{equation}
	\operatorname{Im}\Sigma_{\bm q} = \frac{N_z}{N}\sum_{\bm k}(\bm F_{\bm q}\cdot \bm \pi_{\bm q})^2 {\mathcal V_{q,k}^2}\delta({\varepsilon  -\varepsilon_k - \varepsilon_{q-k}}),
\end{equation}
where $\sum_{\bm k}\delta({\varepsilon  -\varepsilon_k - \varepsilon_{q-k}})$ is actually the two-magnon density of states (DOS). The two-magnon DOS is only non-zero between $2\Delta$ and $2\Delta + 2Jz$, and exhibits a step at the lower end of the continuum. Therefore, there is only a finite plasmon lifetime broadening for plasmon frequencies between these two bounds. 

\section{Hematene}
Having established that plasmon lifetimes can be used to probe spin excitations in antiferromagnets, we now turn to a more realistic scenario. We consider the newly synthesized two-dimensional oxide Hematene, which is a 2D atomic layer of iron oxide (also known as hematite), \ch{Fe2O3}. Hematene is expected to be antiferromagnetic, but the exact nature of the spin excitations remains the subject of some debate. It is therefore an excellent candidate for the use of plasmon lifetime probing.

We show the crystal structure of Hematenen in ???. The \ch{Fe} atoms form a hexagonal lattice, connect by the \ch{O} atoms. We use the suggest antiferromagnetic ordering by \cite{gonzalez}, where the two \ch{Fe} atoms at the unit cell center are ferromagnetically aligned, as well as the other two \ch{Fe} atoms. In addition, we use their calculated exchange interactions, which are all antiferromagnetic. We summarize the crystal structure and exchange Hamiltonian in App. ???. We linearize the exchange Hamiltonian assuming the antiferromagnetic ordering, to obtain the spin-wave dispersion as shown in Fig.???. 

Importantly, the \ch{O} atoms break the bond inversion symmetry, thus allowing for a finite polarization. As outline in Sec.???, we determine the polarization vectors by forming loops of three atoms---in this case of two \ch{Fe} atoms and one \ch{O} atom. We give the resulting polarization vectors between nearest neighbor \ch{Fe} atoms in App. ???. 

Following the derivation as performed for the square antiferromagnet, we can now directly write down the plasmon self-energy as
\begin{equation}
	\Sigma_{\bm q} = \frac{N_z}{N}\sum_{\bm k}\sum_{ij}(\bm F_{\bm q}\cdot \bm \pi^{i,j}_{\bm q})^2 \frac{\mathcal |\mathcal V_{q,k}^{i,j}|^2}{\varepsilon + \ii 0^+ -\varepsilon_k^i - \varepsilon_{q-k}^j},
\end{equation}
where $i,j$ label the four \ch{Fe} atoms in the unit cell. Here, $\mathcal V_{q,k}\equiv \mathcal T^\dagger V_{q,k}\mathcal T$, with $V_{q,k}$ a $4\times4$ matrix given in App.??? and $\mathcal T$ a paraunitary matrix which diagonalizes the antiferromagnetic spin-wave Hamiltonian. 

We can now evaluate the self-energy numerically, by integrating over the Brillioun zone. We show the resulting plasmon lifetime broadening in Fig.???



